\section*{Capítulo \ref{ch:g2:measure} --- Dinero y medidas}

\begin{solution}{exo:g2:measure:1}
Lápiz: cm. Largo de la clase: m. Pluma: g. Mochila: kg. Botella de
agua: L.
\end{solution}

\begin{solution}{exo:g2:measure:2}
$2$ m $= 200$ cm; \quad $300$ cm $= 3$ m; \quad $1$ kg $= 1000$ g.
\end{solution}

\begin{solution}{exo:g2:measure:3}
Por ejemplo: $16 = 10 + 5 + 1$; \ $23 = 20 + 2 + 1$; \
$35 = 20 + 10 + 5$. (Hay otras maneras.)
\end{solution}

\begin{solution}{exo:g2:measure:4}
$10 - 6 = 4$; \quad $20 - 17 = 3$; \quad $50 - 32 = 18$.
\end{solution}

\begin{solution}{exo:g2:measure:5}
Aguja larga en el $6$ y corta entre el $4$ y el $5$: las cuatro y media
($4{:}30$). Aguja larga en el $3$ y corta justo pasado el $8$: las ocho y
cuarto ($8{:}15$).
\end{solution}

\begin{solution}{exo:g2:measure:6}
Las dos y media: la aguja larga apunta abajo, al $6$, y la aguja corta
queda entre el $2$ y el $3$.
\end{solution}

\begin{solution}{exo:g2:measure:7}
Después de mayo: junio. Antes de enero: diciembre. (El mes del cumpleaños varía.)
\end{solution}

\begin{solution}{exo:g2:measure:8}
Gasta: $3 + 4 = 7$. Cambio: $10 - 7 = 3$.
\end{solution}

\begin{solution}{exo:g2:measure:9}
$1$ kg $= 1000$ g, que es más que $800$ g: pesa más el azúcar,
$1000 - 800 = 200$ g más.
\end{solution}

\begin{solution}{exo:g2:measure:10}
Por ejemplo, $20 + 10 + 5 + 5 = 40$ y $10 + 10 + 10 + 10 = 40$
(dos maneras con exactamente cuatro monedas o billetes).
\end{solution}
